%% =====================================================================
%%  第 0 章　阅读须知
%%  源文件：00-阅读须知.md
%%  说明：原文件里的「文档信息」表已提到前置页（main.tex），此处不再重复；
%%        「变更记录」按要求移除。
%% =====================================================================
\chapter{阅读须知}

\begin{guideblock}
\textbf{本章解决什么问题}：帮助你快速建立对本手册的整体认知——它讲什么、写给谁、怎么用、边界在哪里。
\end{guideblock}

值得注意的是，“如何跟人说话”这件事——一个在人类历史上从未需要被专门教学的问题——如今已经发展成了一门需要系统学习的技能。这背后反映了一个不太容易被察觉的趋势，本手册将在后续章节中逐步展开。

先给结论：\textbf{本手册不是一本讲情商、讲话术、讲“高情商回复”的书，而是一份面向实操的接口文档。}它假设你已经在与机器的通信中形成了高效的表达习惯，而你现在需要把这套习惯，迁移到一个更古老、更模糊、也更难预测的对象上——人。

\hcirule

\section{什么是“人类接口”}

\textbf{人类接口}，指的是两个人类之间，通过语言、表情、停顿、沉默等多种手段，实现信息交换的一整套机制。

它具有以下几个基本特征：

\begin{itemize}[leftmargin=2.4em, itemsep=0.3em, topsep=0.3em, parsep=0.2em]
\item \textbf{历史悠久}——早于机器通信出现，一度是人类唯一的通信方式；
\item \textbf{覆盖面广}——在机器通信高度普及的今天，它依然是所有涉及人际交往的场景中的默认通道；
\item \textbf{缺乏规范}——与机器通信不同，它没有统一的接口定义，也没有可供查阅的协议文档。
\end{itemize}

换句话说，人类接口是一个\textbf{运行了很久、却从未被正式定义过}的系统。

\section{为什么需要一本手册}

值得注意的是，机器与机器之间的通信，以下四点是默认成立的：

\begin{manstable}{P{2.4cm} L L}
\tbh{维度} & \tbh{机器之间} & \tbh{人类之间} \\
\midrule
\textbf{格式约定} & 有，且写在双方都能读到的文档里 & 无。格式靠各自经验归纳，两份归纳经常对不上 \\
\textbf{错误码} & 有，失败会明确返回失败 & 无。失败表现为沉默、拖延，或一句“没事” \\
\textbf{重传} & 失败自动重传 & 不重传。没听懂的一方默认装作听懂了 \\
\textbf{载荷完整性} & 收到什么就是什么 & 只收到一部分，剩下的藏在语气和停顿里 \\
\bottomrule
\end{manstable}

这\textbf{不是一份缺陷清单}。需要强调的是，这套机制在它诞生的环境里运行得相当良好——那时候一个人一生要打交道的人不超过两百个，格式可以靠长年累月共同生活慢慢对齐，谁是什么脾气，不需要说。

\textbf{但环境已经变了。}这不仅仅是一次通信方式的迭代，更是一次人际协作底层逻辑的重构。

综上所述，人类接口的“不可靠”，不是它坏了，而是它的使用场景从“熟人社会”扩展到了“陌生人社会”，而它本身并没有随之升级。

\section{读之前，请先接受三条前提}

\hciTag{注意事项}：以下三条属于本手册的适用边界，若不能接受，后续内容可能会出现理解偏差。

\textbf{一、人类接口没有固定格式。}本手册给出的所有格式，都是经验归纳，不是规范，随时可能不成立。这取决于具体的对象、场合与关系状态。

\textbf{二、说得越轻，事越大。}对方把话讲得越轻描淡写，问题往往越严重。这一条贯穿全书，请务必记住。

\textbf{三、本手册改变不了对方。}它只改变你发送和接收的方式。对方回不回你、多快回你、用什么语气回你，不在本手册的调节范围内。

\begin{alertbox}[重要提示]
本手册的部分内容由归纳推理生成，可能包含幻觉（幻觉：机器生成看似合理、实际并不存在的内容）。在应用到真实人际场景之前，请与你身边的真人核实。
\end{alertbox}

\section{如何使用本手册}

本手册采用\textbf{自学体例}编排。每一章均包含：适用场景、分节正文、常见错误、本章小结、练习、自测题与参考答案。

关于练习，只有一条硬性要求：

\begin{mdquote}
\textbf{练习必须对真人执行。不接受纸面推演，不接受“我知道该怎么做了”。}
\end{mdquote}

之所以如此规定，原因在于：人说话时的真实反馈，只能从真人身上取得，本手册无法为你模拟一个活人。此外，练习做错并不扣分；但若跳过不做，后续章节可能会读不懂——这不是危言耸听，而是由本手册的知识结构决定的。

\section{阅读顺序}

请从第 1 章开始，顺序不要调换。

值得注意的是，前十二章的排列并非随意，而是遵循了一条清晰的递进逻辑：\textbf{先解决“你说的话对方没接到”，再解决“对方说的话你没接到”，最后才处理“两边同时说、谁都接不到”。}跳过前面的直接读后面，你将得到一套无法落地执行的方法论。

至于为什么必须遵循这一顺序，第 13 章会给出答案。

\section{一条附加说明}

这本手册第一部分到第三部分教的全部技能，人在三岁左右就已经会了。

这本手册之所以需要存在，是因为这套本事正在丢掉。

请翻到第 1 章。
